\documentclass[11pt]{article}
\usepackage[margin=1in]{geometry}
\usepackage{amsmath, amssymb, amsthm, enumitem}

\title{\textbf{GRE Mathematics Subject Test}\\[4pt] Abstract Algebra Practice Problems -- Set 2}
\author{}
\date{}

\newcounter{prob}
\newcommand{\problem}{\stepcounter{prob}\medskip\noindent\textbf{Problem \theprob.}\ }

\begin{document}
\maketitle
\thispagestyle{empty}

%======================================================================
\section*{A. Groups: Basics \& Subgroups}
%======================================================================

\problem
Let $G$ be a group with $|G|=8$. Which of the following cannot be the order of an element of $G$?
\begin{enumerate}[label=(\alph*)]
\item 3
\item 4
\item 8
\item 2
\end{enumerate}

\problem
Let $G = \mathbb{Z}/2\mathbb{Z}\times\mathbb{Z}/2\mathbb{Z}\times\mathbb{Z}/2\mathbb{Z}$. What is the maximum order of an element in $G$?

\problem
How many subgroups does $S_3$ have?

\problem
True or False: If every element of a group $G$ satisfies $g^2 = e$, then $G$ is abelian.

\problem
In the dihedral group $D_4$ (order 8, symmetries of a square), what is $|Z(D_4)|$?

\problem
Let $G$ be a finite group and $a,b\in G$ with $\operatorname{ord}(a)=3$ and $\operatorname{ord}(b)=5$. If $G$ is abelian, what is $\operatorname{ord}(ab)$?

%======================================================================
\section*{B. Cyclic Groups, Cosets \& Lagrange}
%======================================================================

\problem
True or False: $\mathbb{Z}/4\mathbb{Z}\times\mathbb{Z}/6\mathbb{Z}$ is cyclic.

\problem
Let $H=3\mathbb{Z}$ be a subgroup of $(\mathbb{Z},+)$. List the distinct left cosets of $H$ in $\mathbb{Z}$.

\problem
In $\mathbb{Z}/24\mathbb{Z}$, how many elements have order exactly 8?

\problem
Let $G$ be a cyclic group of order 30. How many subgroups does $G$ have?

\problem
Let $G$ be a group with $|G|=p^2$ where $p$ is prime. Which of the following must be true? Select \textbf{all} that apply.
\begin{enumerate}[label=(\alph*)]
\item $G$ is abelian.
\item $G$ is cyclic.
\item $Z(G)\neq\{e\}$.
\end{enumerate}

%======================================================================
\section*{C. Normal Subgroups, Quotients \& Homomorphisms}
%======================================================================

\problem
Let $\phi\colon G\to H$ be a surjective group homomorphism with $|G|=24$ and $|\ker\phi|=4$. What is $|H|$?

\problem
Is $\{e, (1\;2)(3\;4), (1\;3)(2\;4), (1\;4)(2\;3)\}$ a normal subgroup of $S_4$?

\problem
How many group homomorphisms are there from $\mathbb{Z}$ to $\mathbb{Z}/5\mathbb{Z}$?

\problem
Let $G=\mathbb{Z}/4\mathbb{Z}\times\mathbb{Z}/2\mathbb{Z}$ and let $N=\langle(2,0)\rangle$. What is $G/N$ isomorphic to?

\problem
Which of the following quotient groups is isomorphic to $\mathbb{Z}/3\mathbb{Z}$?
\begin{enumerate}[label=(\alph*)]
\item $\mathbb{Z}/9\mathbb{Z}\,/\,\langle 3\rangle$
\item $\mathbb{Z}/6\mathbb{Z}\,/\,\langle 2\rangle$
\item $S_3/A_3$
\end{enumerate}

%======================================================================
\section*{D. Permutation Groups \& Sylow Theorems}
%======================================================================

\problem
How many permutations in $S_5$ have cycle type $(2,2,1)$ (i.e., are products of two disjoint transpositions)?

\problem
What is the order of the permutation $(1\;2)(3\;4\;5\;6)(7\;8\;9)$ in $S_9$?

\problem
In $S_4$, how many elements are there in the alternating group $A_4$?

\problem
Let $|G|=28=2^2\cdot 7$. Show that $G$ has a normal subgroup of order 7.

\problem
Let $|G|=200=2^3\cdot 5^2$. What are the possible values of $n_5$ (the number of Sylow 5-subgroups)?

%======================================================================
\section*{E. Rings \& Ideals}
%======================================================================

\problem
In $\mathbb{Z}/15\mathbb{Z}$, find all units (invertible elements).

\problem
True or False: Every finite integral domain is a field.

\problem
Let $F$ be a field and $I$ an ideal of $F$. What are the possible ideals?

\problem
Is $\mathbb{Q}[x]/(x^2-2)$ a field? If so, describe it.

\problem
Which of the following rings are PIDs? Select \textbf{all} that apply.
\begin{enumerate}[label=(\alph*)]
\item $\mathbb{Z}$
\item $\mathbb{Q}[x]$
\item $\mathbb{Z}[x]$
\item $\mathbb{Z}[x,y]$
\end{enumerate}

%======================================================================
\section*{F. Fields, Polynomials \& Miscellaneous}
%======================================================================

\problem
Is $x^4 + 1$ irreducible over $\mathbb{Q}$?

\problem
How many elements does the multiplicative group $\mathbb{F}_{16}^*$ have? Is it cyclic?

\problem
Find the degree $[\mathbb{Q}(\sqrt{2},\sqrt{3}):\mathbb{Q}]$.

\problem
How many abelian groups of order $72 = 2^3\cdot 3^2$ are there, up to isomorphism?


%======================================================================
\newpage
\section*{Answer Key}
%======================================================================

\begin{enumerate}

\item (a). By Lagrange's theorem, the order of any element must divide $|G|=8$. The divisors of 8 are 1, 2, 4, 8. Since 3 does not divide 8, no element can have order 3.

\item 2. Every non-identity element $(a,b,c)$ satisfies $(a,b,c)+(a,b,c)=(0,0,0)$ since each coordinate is in $\mathbb{Z}/2\mathbb{Z}$. So every non-identity element has order 2.

\item 6 subgroups. $\{e\}$; three subgroups of order 2: $\{e,(1\;2)\}$, $\{e,(1\;3)\}$, $\{e,(2\;3)\}$; one subgroup of order 3: $A_3=\{e,(1\;2\;3),(1\;3\;2)\}$; and $S_3$ itself.

\item True. If $g^2=e$ for all $g$, then $g=g^{-1}$ for all $g$. For any $a,b\in G$: $ab=(ab)^{-1}=b^{-1}a^{-1}=ba$.

\item $|Z(D_4)|=2$. The center of $D_4$ is $\{e, r^2\}$ where $r$ is a $90^\circ$ rotation. Only the identity and the $180^\circ$ rotation commute with every element.

\item $\operatorname{ord}(ab)=15$. Since $G$ is abelian, $(ab)^{15}=a^{15}b^{15}=e$. Also $\operatorname{ord}(ab)$ must be divisible by both 3 and 5 (since $\gcd(3,5)=1$), so $\operatorname{ord}(ab)=\operatorname{lcm}(3,5)=15$.

\item False. $\mathbb{Z}/4\mathbb{Z}\times\mathbb{Z}/6\mathbb{Z}$ is cyclic iff $\gcd(4,6)=1$, but $\gcd(4,6)=2\neq 1$.

\item There are 3 distinct cosets: $0+3\mathbb{Z}=3\mathbb{Z}=\{\ldots,-3,0,3,6,\ldots\}$, $1+3\mathbb{Z}=\{\ldots,-2,1,4,7,\ldots\}$, and $2+3\mathbb{Z}=\{\ldots,-1,2,5,8,\ldots\}$.

\item An element $\bar{k}$ has order 8 in $\mathbb{Z}/24\mathbb{Z}$ iff $24/\gcd(k,24)=8$, i.e., $\gcd(k,24)=3$. The values of $k\in\{0,\ldots,23\}$ with $\gcd(k,24)=3$ are: $k=3,9,15,21$. So there are $\varphi(8)=4$ elements of order 8.

\item The subgroups of $\mathbb{Z}/30\mathbb{Z}$ are in bijection with divisors of 30. $30=2\cdot 3\cdot 5$, so the number of divisors is $(1+1)(1+1)(1+1)=8$ subgroups.

\item (a) and (c). Groups of order $p^2$ are abelian (this follows from the class equation and the fact that $Z(G)\neq\{e\}$ for $p$-groups). So (a) and (c) are true. (b) is false: $(\mathbb{Z}/p\mathbb{Z})^2$ has order $p^2$ but is not cyclic.

\item $|H|=|G|/|\ker\phi|=24/4=6$ (by the First Isomorphism Theorem, $G/\ker\phi\cong\operatorname{im}\phi=H$).

\item Yes. This is the Klein four-group $V_4=\{e,(1\;2)(3\;4),(1\;3)(2\;4),(1\;4)(2\;3)\}$, and it is normal in $S_4$. One can verify that it is closed under conjugation by any element of $S_4$, or note that it is the kernel of the homomorphism $S_4\to S_3$ given by the action on the three partitions of $\{1,2,3,4\}$ into pairs.

\item There are 5 homomorphisms. A homomorphism $\phi\colon\mathbb{Z}\to\mathbb{Z}/5\mathbb{Z}$ is determined by $\phi(1)$, which can be any element of $\mathbb{Z}/5\mathbb{Z}$ (there are no constraints since $\mathbb{Z}$ is free). So there are 5 homomorphisms.

\item $N=\{(0,0),(2,0)\}$, so $|N|=2$ and $|G/N|=8/2=4$. $G/N$ is abelian of order 4. The coset $(1,0)+N$ has order 2, since $2(1,0)=(2,0)\in N$. Similarly $(0,1)+N$ has order 2, since $2(0,1)=(0,0)\in N$. And $(1,1)+N$ has order 2, since $2(1,1)=(2,0)\in N$. Every non-identity element of $G/N$ has order 2, so $G/N\cong V_4\cong\mathbb{Z}/2\mathbb{Z}\times\mathbb{Z}/2\mathbb{Z}$.

\item (a). $\langle 3\rangle=\{0,3,6\}$ in $\mathbb{Z}/9\mathbb{Z}$, so $\mathbb{Z}/9\mathbb{Z}/\langle 3\rangle$ has order $9/3=3$ and is cyclic, so $\cong\mathbb{Z}/3\mathbb{Z}$. For (b): $\langle 2\rangle=\{0,2,4\}$ in $\mathbb{Z}/6\mathbb{Z}$, so $\mathbb{Z}/6\mathbb{Z}/\langle 2\rangle$ has order $6/3=2$, so $\cong\mathbb{Z}/2\mathbb{Z}$. For (c): $S_3/A_3$ has order $6/3=2$, so $\cong\mathbb{Z}/2\mathbb{Z}$. Only \textbf{(a)}.

\item $\binom{5}{2}\cdot\frac{1}{2}\binom{3}{2}=10\cdot\frac{3}{2}$... More carefully: choose the first transposition by picking 2 elements from 5: $\binom{5}{2}=10$ ways. Choose the second transposition from the remaining 3 elements: $\binom{3}{2}=3$ ways. Divide by $2!$ since the order of the two transpositions doesn't matter: $10\cdot 3/2=15$.

\item $\operatorname{ord}=\operatorname{lcm}(2,4,3)=12$.

\item $|A_4|=|S_4|/2=24/2=12$.

\item $n_7\mid 4$ and $n_7\equiv 1\pmod{7}$. Divisors of 4: 1, 2, 4. Among these, only $1\equiv 1\pmod{7}$. So $n_7=1$, meaning there is a unique (hence normal) Sylow 7-subgroup of order 7.

\item $n_5\mid 8$ and $n_5\equiv 1\pmod{5}$. Divisors of 8: 1, 2, 4, 8. Modulo 5: $1\equiv 1$, $2\equiv 2$, $4\equiv 4$, $8\equiv 3$. Only $n_5=1$ satisfies both conditions. So $n_5=1$.

\item The units are the elements coprime to 15: $\bar{1},\bar{2},\bar{4},\bar{7},\bar{8},\bar{11},\bar{13},\bar{14}$. There are $\varphi(15)=(1-\frac{1}{3})(1-\frac{1}{5})\cdot 15=8$ units.

\item True. This is a standard theorem: in a finite integral domain, every non-zero element $a$ gives an injective map $x\mapsto ax$ (since no zero divisors), and injectivity on a finite set implies surjectivity, so $1$ is in the image, meaning $a$ has an inverse.

\item The only ideals of a field $F$ are $(0)=\{0\}$ and $F$ itself. (If $I$ contains any non-zero element $a$, then $a^{-1}\cdot a = 1\in I$, so $I=F$.)

\item Yes, $\mathbb{Q}[x]/(x^2-2)$ is a field. Since $x^2-2$ is irreducible over $\mathbb{Q}$ (it has no rational roots), the ideal $(x^2-2)$ is maximal in $\mathbb{Q}[x]$ (because $\mathbb{Q}[x]$ is a PID), so the quotient is a field. It is isomorphic to $\mathbb{Q}(\sqrt{2})=\{a+b\sqrt{2}:a,b\in\mathbb{Q}\}$.

\item (a) and (b). $\mathbb{Z}$ is a PID (every ideal is of the form $n\mathbb{Z}$). $\mathbb{Q}[x]$ is a PID ($F[x]$ is a PID when $F$ is a field). $\mathbb{Z}[x]$ is NOT a PID (e.g., $(2,x)$ is not principal). $\mathbb{Z}[x,y]$ is not a PID either (e.g., $(x,y)$ is not principal).

\item Yes, $x^4+1$ is irreducible over $\mathbb{Q}$. Apply Eisenstein's criterion to $f(x+1)=(x+1)^4+1=x^4+4x^3+6x^2+4x+2$: with $p=2$, we have $2\mid 4, 2\mid 6, 2\mid 4, 2\mid 2$, but $4=2^2\nmid 2$ (the constant term). So Eisenstein applies and $f(x+1)$ is irreducible, hence $f(x)$ is irreducible over $\mathbb{Q}$.

\item $|\mathbb{F}_{16}^*|=16-1=15$. Yes, the multiplicative group of any finite field is cyclic. So $\mathbb{F}_{16}^*\cong\mathbb{Z}/15\mathbb{Z}$.

\item $[\mathbb{Q}(\sqrt{2},\sqrt{3}):\mathbb{Q}]=4$. We have $[\mathbb{Q}(\sqrt{2}):\mathbb{Q}]=2$ (minimal polynomial $x^2-2$). Then $\sqrt{3}\notin\mathbb{Q}(\sqrt{2})$ (since $\sqrt{3}=a+b\sqrt{2}$ leads to $3=a^2+2b^2+2ab\sqrt{2}$, forcing $ab=0$, which gives a contradiction). So $[\mathbb{Q}(\sqrt{2},\sqrt{3}):\mathbb{Q}(\sqrt{2})]=2$, and by the tower law, $[\mathbb{Q}(\sqrt{2},\sqrt{3}):\mathbb{Q}]=2\cdot 2=4$.

\item $72=2^3\cdot 3^2$. For $2^3$: partitions of 3 give $\mathbb{Z}/8\mathbb{Z}$, $\mathbb{Z}/4\mathbb{Z}\times\mathbb{Z}/2\mathbb{Z}$, $(\mathbb{Z}/2\mathbb{Z})^3$ (3 choices). For $3^2$: $\mathbb{Z}/9\mathbb{Z}$ or $(\mathbb{Z}/3\mathbb{Z})^2$ (2 choices). Total: $3\times 2 = \mathbf{6}$ abelian groups.

\end{enumerate}

\end{document}
